% Inbuilt themes in beamer

\documentclass[aspectratio=169]{beamer}
\usepackage[export]{adjustbox}
\usepackage{xcolor}

\definecolor{darkred}{RGB}{176, 31, 36}
\definecolor{darkgreen}{RGB}{44, 95, 47}

\usepackage{listings}
\usepackage{booktabs}
\usepackage{graphicx}

% Title page details: 
\title{Tema 3.4: Transformers}
\subtitle{Arquitectura y aplicaciones}

% Theme choice:
\usetheme{Madrid}
\usecolortheme{seagull}
\useinnertheme{circles}

\author{Pablo Pérez Núñez}
 
\setbeamertemplate{navigation symbols}{}
\setbeamertemplate{author}{}

% Ajusta las dimensiones de la diapositiva

\setbeamertemplate{blocks}[rounded][shadow=false]
% \setbeamertemplate{blocks}[default]
\definecolor{myblue}{RGB}{27, 122, 194}
\definecolor{lightblue}{RGB}{55, 150, 225}
\definecolor{lightbluetwo}{RGB}{164, 209, 242}

\setbeamercolor{block title}{bg=lightblue,fg=white}
\setbeamercolor{block body}{bg=lightbluetwo,fg=black}

\setbeamercolor{title}{fg=white, bg =lightblue}
\setbeamercolor{frametitle}{fg=white, bg = lightblue}

\setbeamercolor*{palette secondary}{fg=white,bg=lightblue}

\defbeamertemplate*{footline}{infolines}
{
  \leavevmode%
  \hbox{%
  \begin{beamercolorbox}[wd=\paperwidth,ht=2.25ex,dp=1ex,left]{title in head/foot}%
    \hspace*{2ex} \usebeamerfont{title in head/foot}\insertshorttitle \ - \insertsubtitle
  \end{beamercolorbox}
  % \begin{beamercolorbox}[wd=.25\paperwidth,ht=2.25ex,dp=1ex,right]{date in head/foot}%
  %   \usebeamerfont{date in head/foot}\hspace*{2em}
  %   \insertframenumber{} / \inserttotalframenumber\hspace*{2ex} 
  %\end{beamercolorbox}
  }%
  \vskip0pt%
}

\setbeamertemplate{title page}{
  \centering
  % \usefonttheme{serif}
  \begin{minipage}{0.95\textwidth} % Ajustar el ancho según sea necesario

    \vspace{1em}
    {\color{myblue}\rule{\linewidth}{2pt}\par}
    \vspace{.5cm}
    {\bfseries \huge \textcolor{myblue}{\inserttitle}\par }
    \vspace{1em}
    {\bfseries \LARGE \textcolor{myblue}{\insertsubtitle}\par }
      
    \vspace{2.5em}

    \begin{columns}
      \column{0.6\textwidth}
      \vspace{6em}\\
      {\bfseries \large \textcolor{myblue}{Aprendizaje profundo / Deep Learning}} \\
      \vspace{1em}
      {\textcolor{myblue}{\insertauthor}}
      \column{0.4\textwidth}
      \vfill
      \centering
      \includegraphics[height=4.5em]{imgs/LogoAEPIA.pdf} \\
      \vspace{1em}
      \includegraphics[height=2.5em]{imgs/UIMP.png}
    \end{columns}

    \vspace{1em}

    {\color{myblue}\rule{\linewidth}{.5pt}\par}

  \end{minipage}
    

  %  \usefonttheme{default}
}

% Define the template for section slides
\AtBeginSection[]{
  \begin{frame}
    \vfill
    \LARGE\insertsectionhead\par
    {\color{myblue}\rule{\linewidth}{1pt}\par}
    \vfill
  \end{frame}
}
 

% MIS COSAS --------------------------------
% Bloque con margen inferior
\newenvironment{blockm}[1]{%
  \begin{block}{\textbf{#1}}%
  }{%
  \end{block}%
  \vspace{1em}%
}
% Dos imágenes
\newcommand{\twoimages}[3]{%
  \begin{columns}
    \begin{column}{0.5\textwidth}
      \centering
      \includegraphics[width=#3\textwidth]{#1}
    \end{column}
    \begin{column}{0.5\textwidth}
      \centering
      \includegraphics[width=#3\textwidth]{#2}
    \end{column}
  \end{columns}
}

% Comando para mostrar el título de la subsección al inicio de cada nueva subsección
\AtBeginSubsection[]
{
  \begin{frame}
    \vfill
    \small\insertsectionhead\par
    {\color{myblue}\rule{\linewidth}{1pt}\par}
	\LARGE\insertsubsectionhead\par
    \vfill
  \end{frame}
}

% ------------------------------------------

\begin{document}

% Title page frame
\begin{frame}[plain]
	\titlepage 
\end{frame}

\logo{}


\subsection{Transformers: Arquitectura}

\begin{frame}{Transformers}
  
  Una vez conocido el funcionamiento de los \textbf{mecanismos de atención}, podemos empezar a hablar del \textbf{Transformer}.\\
  \vspace{.5cm}
  \begin{block}{}
    Presentado en 2017 por Vaswani et al. Fué diseñado para tareas de traducción (seq-2-seq), pero hoy en día es clave en la mayoría de tareas NLP.
  \end{block}
  \vspace{.5cm}
  \textbf{El Transformer:}
  \begin{itemize}
    \item Es una arquitectura (conjunto de capas).
    \item Utiliza mecanismos de self y cross-attention.
    \item Es altamente paralelizable: permite entrenar modelos más grandes en menos tiempo.
  \end{itemize}
  % Un transformador no es sólo una capa de autoatención, es una arquitectura. No está muy claro qué se considera un transformador y qué no, pero aquí utilizaremos la siguiente definición:
  % Cualquier arquitectura diseñada para procesar un conjunto conectado de unidades -como los tokens de una secuencia o los píxeles de una imagen- en la que la única interacción entre unidades se produce a través de la autoatención.
  % Al igual que con otros mecanismos, como las convoluciones, ha surgido un enfoque más o menos estándar sobre cómo construir capas de autoatención en una red más grande. El primer paso es envolver la autoatención en un bloque que podamos repetir.
\end{frame}

\begin{frame}{Transformers}
  \begin{columns}
    \begin{column}{.4\textwidth}
      \textbf{El Transformer tradicional se compone de dos módulos:}
      \begin{itemize}
        \item Encoder
        \item Decoder
      \end{itemize}
      \vspace{.5cm}
      \textbf{Cada uno está formado por:}
      \begin{itemize}
        \item Capas self-attention.
        \item Capas cross-attention.
        \item Capas de normalización.
        \item Perceptrones multicapa.
        \item Conexiones residuales.
      \end{itemize}
    \end{column}
    \begin{column}{.6\textwidth}
      \includegraphics[width=\textwidth, center]{imgs/tema4/att/Transformer_all.pdf}
    \end{column}
  \end{columns}
\end{frame}

\begin{frame}{Transformers}
  \begin{columns}
    \begin{column}{.4\textwidth}
      \textbf{El encoder:}
      \begin{itemize}
        \item Recibe secuencia de entrada.
        \item Genera secuencia codificada.
      \end{itemize}
      \vspace{.5cm}
      \textbf{El decoder:}
      \begin{itemize}
        \item Recibe salida del encoder.
        \item Recibe su salida anterior.
        \begin{itemize}
          \item Primero recibe \emph{$<start>$}.
        \end{itemize}
        \item Genera hasta \emph{$<end>$}.
      \end{itemize}
    \end{column}
    \begin{column}{.6\textwidth}
      \includegraphics[width=\textwidth, center]{imgs/tema4/att/Transformer_all_s1.pdf}
    \end{column}
  \end{columns}
\end{frame}


\begin{frame}{Transformers}
  \begin{columns}
    \begin{column}{.4\textwidth}
      \textbf{El encoder:}
      \begin{itemize}
        \item Recibe secuencia de entrada.
        \item Genera secuencia codificada.
      \end{itemize}
      \vspace{.5cm}
      \textbf{El decoder:}
      \begin{itemize}
        \item Recibe salida del encoder.
        \item Recibe su salida anterior.
        \begin{itemize}
          \item Primero recibe \emph{$<start>$}.
        \end{itemize}
        \item Genera hasta \emph{$<end>$}.
      \end{itemize}
    \end{column}
    \begin{column}{.6\textwidth}
      \includegraphics[width=\textwidth, center]{imgs/tema4/att/Transformer_all_s2.pdf}
    \end{column}
  \end{columns}
\end{frame}


\begin{frame}{Transformers}
  \begin{columns}
    \begin{column}{.4\textwidth}
      \textbf{El encoder:}
      \begin{itemize}
        \item Recibe secuencia de entrada.
        \item Genera secuencia codificada.
      \end{itemize}
      \vspace{.5cm}
      \textbf{El decoder:}
      \begin{itemize}
        \item Recibe salida del encoder.
        \item Recibe su salida anterior.
        \begin{itemize}
          \item Primero recibe \emph{$<start>$}.
        \end{itemize}
        \item Genera hasta \emph{$<end>$}.
      \end{itemize}
    \end{column}
    \begin{column}{.6\textwidth}
      \includegraphics[width=\textwidth, center]{imgs/tema4/att/Transformer_all_s3.pdf}
    \end{column}
  \end{columns}
\end{frame}


\begin{frame}{Transformers}
  \begin{columns}
    \begin{column}{.4\textwidth}
      \textbf{El encoder:}
      \begin{itemize}
        \item Recibe secuencia de entrada.
        \item Genera secuencia codificada.
      \end{itemize}
      \vspace{.5cm}
      \textbf{El decoder:}
      \begin{itemize}
        \item Recibe salida del encoder.
        \item Recibe su salida anterior.
        \begin{itemize}
          \item Primero recibe \emph{$<start>$}.
        \end{itemize}
        \item Genera hasta \emph{$<end>$}.
      \end{itemize}
    \end{column}
    \begin{column}{.6\textwidth}
      \includegraphics[width=\textwidth, center]{imgs/tema4/att/Transformer_all_s4.pdf}
    \end{column}
  \end{columns}
\end{frame}

\begin{frame}{Transformers}
  \begin{columns}
    \begin{column}{.4\textwidth}
      \textbf{El encoder:}
      \begin{itemize}
        \item Recibe secuencia de entrada.
        \item Genera secuencia codificada.
      \end{itemize}
      \vspace{.5cm}
      \textbf{El decoder:}
      \begin{itemize}
        \item Recibe salida del encoder.
        \item Recibe su salida anterior.
        \begin{itemize}
          \item Primero recibe \emph{$<start>$}.
        \end{itemize}
        \item Genera hasta \emph{$<end>$}.
      \end{itemize}
    \end{column}
    \begin{column}{.6\textwidth}
      \includegraphics[width=\textwidth, center]{imgs/tema4/att/Transformer_all_s5.pdf}
    \end{column}
  \end{columns}
\end{frame}

\begin{frame}{Transformers: Encoder}
  El \textbf{encoder} de un Transformer se compone \textbf{6 módulos} concatenados.\\
  \vspace{.5cm}
  \textbf{Cada módulo contiene, por orden:}
  \begin{enumerate}
    \item Capa multi-head self-attention.
    \item Conexión residual con entrada previa.
    \item Capa de normalización.
    \item Un MLP por cada elemento.
    \item Conexión residual con entrada previa.
    \item Capa de normalización.
  \end{enumerate}

\end{frame}

\begin{frame}{Transformers: Encoder}
  En detalle:\\
  \vspace{.5cm}
  \includegraphics[width=.85\textwidth, center]{imgs/tema4/att/Transformer_enc.pdf}
\end{frame}

\begin{frame}{Transformers: Decoder}
  El \textbf{decoder} de un Transformer se compone \textbf{6 módulos} concatenados.\\
  \vspace{.5cm}
  \textbf{Cada módulo contiene, por orden:}
  \begin{enumerate}
    \item Capa multi-head self-attention.
    \item Conexión residual con entrada previa.
    \item Capa de normalización.
    \item Capa multi-head \textbf{cross-attention (con salida de encoder)}.
    \item Conexión residual con entrada previa.
    \item Capa de normalización.
    \item Un MLP por cada elemento.
    \item Conexión residual con entrada previa.
    \item Capa de normalización.
  \end{enumerate}\end{frame}

\begin{frame}{Transformers: Decoder}
  En detalle:\\
  \vspace{.5cm}
  \includegraphics[width=.85\textwidth, center]{imgs/tema4/att/Transformer_dec.pdf}
\end{frame}

\begin{frame}{Transformers vs RNN}
  \begin{block}{}
    \textbf{RNN:} Como ya habíamos visto, en tareas seq2seq condensan toda la información de la entrada en un único elemento. Esto se traduce en una \textbf{memoria corto-placista}.
  \end{block}
  \includegraphics[width=.55\textwidth, center]{imgs/tema4/att/ATT_vs_RNN.pdf}
  
  \begin{block}{}
    \textbf{Transformers:} Aprovechan toda la secuencia de entrada para generar la salida.
  \end{block}

\end{frame}

\begin{frame}{Transformers vs RNN}
  \begin{block}{}
    \textbf{RNN:} Al depender un elemento del anterior, \textbf{no es posible paralelizar} su ejecución. Esto implica un menor rendimiento.
  \end{block}

  \includegraphics[width=.55\textwidth, center]{imgs/tema4/att/ATT_vs_RNN_2.pdf}
  
  \begin{block}{}
    \textbf{Transformers:} Altamente paralelizables. Grandes modelos entrenados en menos tiempo.
  \end{block}

\end{frame}
 
\begin{frame}{Transformers vs RNN}
  \begin{block}{}
    \textbf{RNN:} \textbf{Tienen} en cuenta el orden de los elementos de la secuencia de entrada.
  \end{block}
  \vspace{.5cm}
  \textbf{Transformers:}
  \begin{itemize}
    \item \textbf{Tal como los hemos descrito hasta este punto}, no tienen en cuenta el orden.
    \item Esto es una \textbf{desventaja}, el orden de los elementos es muy importante.
    \item Necesitamos incorporar algo más: \textbf{Positional encoding}.
  \end{itemize}
\end{frame}

\begin{frame}{Transformers: Positional encoding}
  \textbf{Positional encoding:}
  \begin{itemize}
    \item Sumar un vector $\mathbf{t}_{i}$ a cada elemento $\mathbf{x}_{i}$ de la entrada.
    \item Este vector ha de ser del mismo tamaño $k$ que $\mathbf{x}_{i}$.
  \end{itemize}
  \vspace{.3cm}
  Los valores se obtienen:
  $$
  \begin{aligned}
  \mathbf{t}_{i,2j} &= \sin\left(\frac{i}{10000^{\frac{2j}{k}}}\right) \\\\
  \mathbf{t}_{i,2j+1} &= \cos\left(\frac{i}{10000^{\frac{2j}{k}}}\right).
  \end{aligned}
  $$
  \vspace{.3cm}
  \begin{block}{}
    Como se puede ver, los elementos pares del vector se obtienen de forma diferente a los impares.    
  \end{block}

\end{frame}

% \begin{frame}{Transformers: Positional encoding}
%   Visualización de los vectores $\mathbf{t}_{i}$ en secuencias de tamaño máximo 50 con $k=128$:\\
%   \vspace{.5cm}
%   \includegraphics[width=.7\textwidth, center]{imgs/tema4/att/positional-encoding.png}
% \end{frame}

\begin{frame}{Transformers: Positional encoding}
  Localización dentro del Transformer:\\
  \vspace{.5cm}
  \includegraphics[width=.7\textwidth, center]{imgs/tema4/att/positional_enc.pdf}
\end{frame}

\subsection{Transformers: Aplicaciones}

\begin{frame}[t]{Transformers: Aplicaciones}
  \begin{block}{}
    Diseñados inicialmente para resolver problemas de traducción. En la actualidad se utilizan en muchas más tareas.
  \end{block}

  \textbf{Natural Language Processing:}
  \begin{itemize}
    \item Traducción
    \item Clasificación
    \item Conversación (chat)
  \end{itemize}

  \textbf{Imágenes:}
  \begin{itemize}
    \item Detección de objetos
    \item Super-resolución
  \end{itemize}

  \textbf{Multimodal:}
  \begin{itemize}
    \item Vídeo y texto (por ejemplo, subtítulos automáticos para videos)
    \item Generación de descripciones para imágenes (Image captioning)
  \end{itemize}

\end{frame}

\begin{frame}[t]{Transformers: Aplicaciones}
  \begin{block}{}
    En tareas NLP, gracias a la introducción de los Transformers los modelos de lenguaje han crecido en tamaño, precisión y complejidad.
  \end{block}
  \vspace{.2cm}
  \includegraphics[width=.75\textwidth, center]{imgs/tema4/att/llm_graph.png}
  \vspace{.2cm}
  El modelo ELMo no utiliza transformers, está basado en LSTM bidireccionales.
\end{frame}

\begin{frame}[t]{Transformers: Aplicaciones}
  \textbf{GPT-4 (Generative Pre-trained Transformer 4):}
  \begin{itemize}
    \item Creado por OpenAI.
    \item Modelo de lenguaje con miles de millones de parámetros.
    \item Utilizado en generación de texto coherente, traducción y chatbots avanzados.
  \end{itemize}
  \vspace{1cm}
  \includegraphics[width=\textwidth, center]{imgs/tema4/att/chatgpt-3.pdf}

\end{frame}

\begin{frame}[t]{Transformers: Aplicaciones}
  \textbf{CLIP (Contrastive Language-Image Pretraining):}
  \begin{itemize}
    \item Modelo creado por OpenAI que combina imágenes y texto utilizando transformers.
    \item Capacidad para entender contenido visual y textual, útil en clasificación de imágenes y búsqueda basada en texto.
  \end{itemize}
  \vspace{.5cm}
  \includegraphics[width=.75\textwidth, center]{imgs/tema4/att/clip.png}
\end{frame}

% \begin{frame}[t]{Transformers: Aplicaciones}
%   \textbf{BERT (Bidirectional Encoder Representations from Transformers):}
%   \begin{itemize}
%     \item Creado por Google.
%     \item Es bidireccional, teniendo en cuenta el contexto anterior y posterior de una palabra.
%     \item Utilizado en tareas de procesamiento de lenguaje natural como clasificación, etiquetado de entidades y resolución de preguntas.
%   \end{itemize}
%   \vspace{.5cm}
%   \includegraphics[width=.6\textwidth, center]{imgs/tema4/att/bert.jpg}
% \end{frame}

% \begin{frame}[t]{Transformers: Aplicaciones}
%   \textbf{T5 (Text-to-Text Transfer Transformer):}
%   \begin{itemize}
%     \item Modelo de lenguaje de Google que sigue el enfoque ``texto a texto''.
%     \item Capaz de realizar múltiples tareas de procesamiento de lenguaje natural, como traducción, resumen y generación de respuestas.
%     \item Versátil y eficiente al unificar diferentes tareas bajo una sola arquitectura.
%   \end{itemize}
%   \vspace{.5cm}
%   \includegraphics[width=.6\textwidth, center]{imgs/tema4/att/t5.jpg}
% \end{frame}

\begin{frame}[t]{Referencias}
  \begin{enumerate}
    \item \href{https://learningturtle.github.io/Blog/posts/attention_another_perspective_part2/}{\textbf{Attention? An Other Perspective!, Nikhil Shah}}\\
    \item \href{https://glouppe.github.io/info8010-deep-learning/?p=lecture7.md}{\textbf{Lecture 7: Attention and transformers, Prof. Gilles Louppe}}\\  
    \item \href{https://peterbloem.nl/blog/transformers}{\textbf{Transformers from scratch, Peter Bloem}}
    \item \href{https://jalammar.github.io/illustrated-transformer/}{\textbf{The Illustrated Transformer}} \\
  \end{enumerate}
\end{frame}



\end{document}
